%! Exponential and Logarithmic Equations and Inequalities — Unit 2, Lesson 2.13 — Warm-Up
\documentclass[10pt]{article}
\usepackage{precalculus-article}
\usepackage{precalculus-boxes}
\newcommand{\TallMath}[1]{$\displaystyle #1\rule[-1.4em]{0pt}{3.2em}$}

\begin{document}

\pageheader{Unit 2, Lesson 2.13}{Warm-Up}
\namedateperiod

\vspace{0.15in}

\begin{enumerate}[leftmargin=1.8em, itemsep=16pt]

\item Solve each equation for $x$. Show your reasoning.

\vspace{6pt}
\begin{enumerate}[leftmargin=1.5em, itemsep=10pt, label=\alph*.]
  \item $2^x = 32$

  \vspace{6pt}
  $x = $ \blank{3cm}

  \item $3^x = \dfrac{1}{27}$

  \vspace{6pt}
  $x = $ \blank{3cm}
\end{enumerate}

\item Let $f(x) = \log_2 x$ and $g(x) = 2^x$.

\vspace{6pt}
Evaluate $g(f(8))$. Show each step.

\vspace{6pt}
Step 1: $f(8) = \log_2 8 = $ \blank{2.5cm}

\vspace{6pt}
Step 2: $g(f(8)) = 2^{\,\square} = $ \blank{2.5cm}

\vspace{8pt}
What property of $f$ and $g$ does this result demonstrate?

\writeline

\item For the function $f(x) = 5 \cdot 2^x + 1$, list the two operations applied to the
      basic exponential $2^x$, in order.

\vspace{6pt}
Operation 1: \blank{6cm}

\vspace{6pt}
Operation 2: \blank{6cm}

\end{enumerate}

\end{document}
